\subsection{Анализ рынка и описание условий}

В последние годы контейнеризация и оркестрация приложений стали одним из ключевых направлений развития корпоративной ИТ-инфраструктуры. Развитие микросервисной архитектуры и контейнерных подходов сделало Kubernetes одной из базовых платформ для эксплуатации современных приложений. Применение Kubernetes позволяет ускорять развёртывание сервисов, повышать масштабируемость систем и упрощать сопровождение распределённых приложений \cite{Kubernetes}.

Вместе с тем рост числа контейнерных сред приводит к увеличению требований к их защищённости, поскольку контейнерные технологии имеют связанные с безопасностью риски, а уязвимости и ошибки конфигурации могут возникать как на этапе сборки, так и на этапе эксплуатации \cite{NISTSP800190,NISTSP800204C}. Ошибки конфигурации, несанкционированные действия пользователей и аномальная активность внутри кластера могут привести к нарушению доступности сервисов, утечке данных и сбоям в работе инфраструктуры \cite{NISTSP800190,NISTSP800204C}.

Одним из наиболее информативных источников данных о действиях в Kubernetes-кластере являются audit-логи, содержащие сведения о запросах к API-серверу, выполнении административных операций и взаимодействии различных субъектов с объектами кластера \cite{KubernetesAudit}. Анализ таких журналов позволяет выявлять отклонения от нормального поведения, обнаруживать подозрительные действия и своевременно реагировать на потенциальные инциденты информационной безопасности. В связи с этим разработка программного средства обнаружения аномалий в Kubernetes-кластере на основе анализа audit-логов является актуальной задачей для организаций, использующих контейнерную инфраструктуру в продуктивной среде.

Спрос на подобные решения формируется, прежде всего, со стороны компаний, для которых Kubernetes уже используется или оценивается для внедрения в производственную среду. По данным CNCF Annual Survey 2024, Kubernetes в производственной среде или на стадии внедрения используют 85\% респондентов, а четверть участников опроса сообщили, что почти вся их разработка и поставка выполняется с применением cloud native-подходов \cite{CNCFAnnualSurvey2024}. Это показывает, что Kubernetes остаётся массовой и развивающейся платформой, а задачи мониторинга, аудита и выявления аномалий сохраняют практическую актуальность. Для таких организаций особенно важны инструменты, позволяющие не только контролировать состояние кластера, но и выявлять нетипичные действия администраторов, сервисных учётных записей и пользователей, имеющих доступ к API Kubernetes.

В рамках данной выпускной квалификационной работы рассматривается гипотетический заказчик, заинтересованный во внедрении программного комплекса для анализа audit-логов Kubernetes-кластера. 
Учитывая новизну и сложность разрабатываемого программного средства, а также необходимость интеграции с существующей инфраструктурой заказчика, можно предположить, что общая стоимость проекта составит порядка 3 000 000 рублей. Данная сумма включает затраты на разработку, тестирование, внедрение и первичную настройку системы. Предполагаемый срок реализации проекта составляет 5–6 месяцев, что соответствует типичным срокам разработки программных решений подобного класса. После завершения разработки планируется тиражирование решения и его адаптация под требования других клиентов. Поскольку Kubernetes широко используется в различных отраслях, включая финансовые и IT-компании \cite{CNCFAnnualSurvey2024}, можно предположить наличие устойчивого спроса на подобные решения. Предположим, что в течение года воспользуются услугой внедрения программного комплекса в среднем 3 организации. Перед очередным внедрением потребуется адаптация комплекса под требования заказчика.